\section{Technology \& Implementation Plan}
\label{sec:technology_implementation}

\subsection{Technology Stack}
\label{subsec:tech_stack}

The GridSense AI prototype is built on a modern, high-performance, and extensible technology stack:
\begin{itemize}[leftmargin=1.5em, itemsep=0.2em]
    \item \textbf{Frontend UI:} \textbf{React + Tailwind CSS} for a responsive, modular control-room interface.
    \item \textbf{Visualization:} \textbf{Recharts / D3.js} for high-performance interactive time-series curves and uncertainty ribbons.
    \item \textbf{Backend REST API:} \textbf{Python + FastAPI} providing asynchronous endpoints for telemetry ingestion, forecast generation, and scenario evaluation.
    \item \textbf{AI/ML Engine:} \textbf{Python + XGBoost} executing multi-step tabular time-series regression and quantile inference.
    \item \textbf{Decision Engine:} Python-based rule and multi-objective optimization logic.
    \item \textbf{Data Store:} \textbf{PostgreSQL / JSON Cache} managing historical time series, weather feeds, and scenario state.
\end{itemize}

\subsection{System Architecture \& Components}
\label{subsec:system_components}

\Cref{fig:system_architecture} illustrates the modular multi-tier system architecture:

\begin{figure}[htbp]
    \centering
    \begin{tikzpicture}[
        node distance=0.32cm and 0.35cm,
        layerBox/.style={rectangle, rounded corners=3pt, draw=accentTeal!80, fill=accentTeal!10, text width=8.2cm, align=center, font=\scriptsize, inner sep=3pt},
        dataBox/.style={rectangle, rounded corners=3pt, draw=accentBlue!80, fill=accentBlue!10, text width=2.45cm, align=center, font=\tiny, inner sep=3pt},
        dashBox/.style={rectangle, rounded corners=3pt, draw=accentAmber!80, fill=accentAmber!10, text width=8.2cm, align=center, font=\scriptsize\bfseries, inner sep=3.5pt},
        arrow/.style={-{Stealth[scale=0.75]}, thick, color=darkSlate}
    ]
        % Top Inputs
        \node[dataBox] (src1) {\textbf{Weather Data}\\Irradiance, Wind, Temp};
        \node[dataBox, right=0.35cm of src1] (src2) {\textbf{Historical Records}\\Generation Time Series};
        \node[dataBox, right=0.35cm of src2] (src3) {\textbf{Site Metadata}\\Capacity, Coordinates};

        % Ingestion Layer
        \node[layerBox, below=0.40cm of src2] (ingest) {\textbf{Data Ingestion \& Validation Layer}\\Format Validation \textbar{} Imputation \textbar{} Timestamp Alignment};

        % Forecasting Layer
        \node[layerBox, below=of ingest] (forecast) {\textbf{Forecasting Engine}\\Multi-Horizon 24h, 48h, 72h Generation Prediction};

        % Uncertainty & Risk Layer
        \node[layerBox, below=of forecast] (risk) {\textbf{Uncertainty \& Risk Engine}\\Quantile Intervals ($p_{10}, p_{50}, p_{90}$) \textbar{} Surplus/Shortfall Flags};

        % Simulation & Optimization Layer
        \node[layerBox, below=of risk] (optim) {\textbf{Simulation \& Decision Engine}\\What-If Parameter Perturbation \textbar{} Dispatch Optimization};

        % Recommendation Layer
        \node[layerBox, below=of optim] (recom) {\textbf{Recommendation Layer \& Explainable Copilot}\\Actionable Advice \textbar{} Cost \& Carbon Rationale};

        % Operator Dashboard
        \node[dashBox, below=of recom] (dash) {Operator Dashboard \& Human Decision Console};

        % Arrows
        \draw[arrow] (src1) -- (src1 |- ingest.north);
        \draw[arrow] (src2) -- (ingest.north);
        \draw[arrow] (src3) -- (src3 |- ingest.north);

        \draw[arrow] (ingest) -- (forecast);
        \draw[arrow] (forecast) -- (risk);
        \draw[arrow] (risk) -- (optim);
        \draw[arrow] (optim) -- (recom);
        \draw[arrow] (recom) -- (dash);
    \end{tikzpicture}
    \caption{GridSense AI System Architecture: layered progression from data ingestion to operator dashboard.}
    \label{fig:system_architecture}
\end{figure}

\subsection{Implementation Flow \& Prototype Scope}
\label{subsec:implementation_flow}

\textbf{End-to-End Implementation Flow:}
\begin{center}
    \small
    \texttt{Input Data} $\longrightarrow$ \texttt{Preparation} $\longrightarrow$ \texttt{Feature Pipeline} $\longrightarrow$ \texttt{Forecasting} $\longrightarrow$ \texttt{Uncertainty \& Risk} $\longrightarrow$ \texttt{Scenario Sandbox} $\longrightarrow$ \texttt{Decision Optimization} $\longrightarrow$ \texttt{Explainable Recommendation} $\longrightarrow$ \textbf{\color{accentAmber}Dashboard}
\end{center}

The prototype implementation executes across three tightly integrated execution stages:
\begin{enumerate}[leftmargin=1.5em, itemsep=0.3em]
    \item \textbf{Data Pipeline \& Forecasting Service:} Upon receiving raw telemetry and NWP weather feeds, the ingestion service performs timestamp synchronization and linear interpolation for missing values. The feature engineering module constructs time-lag and rolling window metrics before dispatching the feature vector to the trained XGBoost model, which returns expected generation curves along with $p_{10}$ and $p_{90}$ uncertainty bounds for the $24$--$72$ hour horizon.
    \item \textbf{Decision \& What-If Execution Core:} The decision service scans the forecast trajectory against anticipated baseline demand to detect potential deficit or curtailment events. If an operator adjusts a what-if slider (e.g., simulating a sudden $30\%$ irradiance drop or battery outage), the engine recalculates risk metrics and ranks dispatch actions based on cost, carbon intensity, and feasibility.
    \item \textbf{Explainable Dashboard Integration:} The FastAPI backend serializes time-series arrays and structured recommendation payloads to the React interface, rendering dynamic uncertainty ribbons, alert banners, and explainable decision cards for human operator review.
\end{enumerate}

\textbf{Hackathon Prototype Scope (MVP):} The hackathon submission focuses on a self-contained, working end-to-end loop: ingesting real-world solar and wind datasets, predicting generation with confidence intervals across $24$--$72$ hours, detecting imbalance risks, simulating interactive what-if perturbations, and displaying explainable decision recommendations on the operator console.
